\section*{अध्याय \ref{ch:b2:metamorphic-facies} --- कायांतरी अभिक्रियाएँ और संलक्षणियाँ}

\begin{solution}{exo:b2:metamorphic-facies:1}
(क) संपर्क \omterm{def:b2:metamorphic-facies:metamorphism}{कायांतरण}, अंतर्वेधन की ऊष्मा से होने वाला एक स्थानीय ऊष्मीय \omterm{def:b2:metamorphic-facies:metamorphism}{कायांतरण}।
(ख) विस्थापन \omterm{def:b2:metamorphic-facies:metamorphism}{कायांतरण}, एक भ्रंश-मंडल में: कण-आकार में कमी, एक मायलोनाइट।
(ग) महासागर-तल \omterm{def:b2:metamorphic-facies:metamorphism}{कायांतरण}, कटक के निकट गर्म समुद्री जल द्वारा, अपूर्ण
पुनःक्रिस्टलन के साथ। (घ) दफ़न \omterm{def:b2:metamorphic-facies:metamorphism}{कायांतरण}, क्षेत्रीय, विरूपण या
शिस्टता के बिना।
\end{solution}

\begin{solution}{exo:b2:metamorphic-facies:2}
$\Delta_r S = 95.40 - 82.80 = \qty{12.60}{J.K^{-1}}$, $\Delta_r V = 4.986 - 4.415 =
\qty{0.571}{J.bar^{-1}}$; ढलान $12.60/0.571 = \qty{22.1}{bar/K} = \qty{2.21}{MPa/K}$।
\qty{100}{\celsius} में साम्य-दाब \qty{2.21}{kbar} बढ़ता है।
\end{solution}

\begin{solution}{exo:b2:metamorphic-facies:3}
(क) हरितशिस्ट। (ख) नीलशिस्ट। (ग) एम्फ़ीबोलाइट: हॉर्नब्लेंड, और ऐसा प्लैजियोक्लेज़ जिसमें
\qty{17}{\percent} से अधिक एनॉर्थाइट हो। (घ) एक्लोजाइट।
\end{solution}

\begin{solution}{exo:b2:metamorphic-facies:4}
(क) $\rho = M/V$: कायनाइट $162.046/44.15 = \qty{3.670}{g.cm^{-3}}$, ऐंडालुसाइट
$162.046/51.52 = 3.145$, सिलिमेनाइट $162.046/49.86 = 3.250$। (ख) कायनाइट: अधिक
दाब पर गिब्स ऊर्जा का पद $V\,(P - 1)$ सबसे छोटे आयतन के लिए सबसे कम बढ़ता है।
(ग) प्राकृतिक क्रिस्टलों में अन्य तत्त्व होते हैं (उदाहरण के लिए, प्राकृतिक ऐंडालुसाइट और सिलिमेनाइट में फ़ेरिक लोहा
कुछ ऐलुमिनियम का स्थान लेता है), और प्राकृतिक नमूनों पर किए गए मापन
तालिकाओं के मोलर आयतनों से कम परिशुद्ध होते हैं।
\end{solution}

\begin{solution}{exo:b2:metamorphic-facies:5}
(क) $\Delta_r G = 3900 - 298.15\times8.59 = \qty{+1339}{J} > 0$: कायनाइट। (ख)
\qty{1}{bar} पर, $\Delta_r H = T\Delta_r S$: $T = 3900/8.59 = \qty{454}{K} =
\qty{181}{\celsius}$। (ग) $P = (8.59\times673.15 - 3900)/0.737 = 1882/0.737 =
\qty{2554}{bar} = \qty{2.55}{kbar}$, मापे गए 2.54 की तुलना में। (घ) निम्न तापमान पर अभिक्रिया
आगे बढ़ने के लिए बहुत धीमी है: ऐंडालुसाइट मितस्थायी रूप में बचा रहता है, जैसे अधिकांश उच्च-तापमान
खनिज सतह पर बचे रहते हैं।
\end{solution}

\begin{solution}{exo:b2:metamorphic-facies:6}
(क) $\Delta_r H = 7700 + 1000\times(121.39 - 122.29) = \qty{6800}{J}$; $\Delta_r S = 296.44
- 285.02 = \qty{11.42}{J.K^{-1}}$; $\Delta_r G = 6800 - 1000\times11.42 = \qty{-4620}{J}$।
(ख) \qty{298.15}{K} के मानों के साथ: $7700 - 1000\times12.60 = \qty{-4900}{J}$।
(ग) सिलिमेनाइट। सरल विधि की त्रुटि \qty{280}{J} है, लगभग \qty{6}{\percent}; इसे
$\Delta_r V = 0.571$ से भाग देने पर रेखा पर लगभग \qty{490}{bar} आता है, जो एन्थैल्पियों की
अनिश्चितता के प्रभाव से बहुत कम है।
\end{solution}

\begin{solution}{exo:b2:metamorphic-facies:7}
(क) $\Delta_r G = -1127.4 + 1128.5 = \qty{1.1}{kJ}$ और $\Delta_r V = (34.15 -
36.93)/10 = \qty{-0.278}{J.bar^{-1}}$: $P = 1100/0.278 = \qty{3957}{bar} \approx
\qty{4.0}{kbar}$। (ख) $\Delta_r S = 88.0 - 91.7 = -3.7$; ढलान $-3.7/-0.278 =
\qty{13.3}{bar/K}$। (ग) $\Delta_r H = 0$ के साथ, $P = T\Delta_r S/\Delta_r V = 13.31\,T$:
\qty{573.15}{K} पर, \qty{7629}{bar} $= \qty{7.63}{kbar}$। (घ) कवचों का ऐरागोनाइट सतह पर
मितस्थायी है: जीव उसे बनाते हैं, और वह बना रहता है क्योंकि वह
कैल्साइट में केवल धीरे-धीरे बदलता है।
\end{solution}

\begin{solution}{exo:b2:metamorphic-facies:8}
(क) ढलान $-32.4/-1.698 = \qty{19.1}{bar/K}$ है। (ख) $P = (T\Delta_r S - \Delta_r H)/\Delta_r V =
(-32.4\times773.15 + 5000)/(-1.698) = (-25050 + 5000)/(-1.698) = \qty{11808}{bar}
\approx \qty{11.8}{kbar}$। (ग) $\sqrt{2.6^2 + 3.6^2 + 1.0^2} = \qty{4.55}{kJ}$, यह
$\Delta_r H$ पर है, अतः दाब पर $4550/1.698 = \qty{2.7}{kbar}$। (घ) एक्लोजाइट संलक्षणी के
दाबों पर प्लैजियोक्लेज़ का एल्बाइट घटक जेडाइट और क्वार्ट्ज़ की तुलना में
अस्थायी है; सोडियम पाइरॉक्सीन (ओम्फ़ैसाइट) में चला जाता है, और प्लैजियोक्लेज़
लुप्त हो जाता है।
\end{solution}

\begin{solution}{exo:b2:metamorphic-facies:9}
(क) 1 से \qty{2}{kbar} तक: $1000/34 = \qty{29.4}{bar/K}$; 2 से 3 तक: $1000/25 =
\qty{40.0}{bar/K}$। (ख) $\Delta_r S$ स्थिर होने पर $\Delta_r V$ ढलान के व्युत्क्रमानुपाती
है: वह $40.0/29.4 = 1.36$ के गुणक से सिकुड़ा। (ग) मुक्त हुआ जल: दाब बढ़ने पर उसका मोलर
आयतन घटता है, जबकि ठोसों के आयतन मुश्किल से बदलते हैं।
\end{solution}

\begin{solution}{exo:b2:metamorphic-facies:10}
(क) \qty{873.15}{K} पर: कायनाइट $=$ ऐंडालुसाइट $(8.59\times873.15 - 3900)/0.737 =
\qty{4885}{bar}$; कायनाइट $=$ सिलिमेनाइट $(12.60\times873.15 - 7700)/0.571 =
\qty{5782}{bar}$; ऐंडालुसाइट $=$ सिलिमेनाइट $(4.01\times873.15 - 3800)/(-0.166) =
\qty{1799}{bar}$। (ख) \qty{4885}{bar} पर, कायनाइट $=$ सिलिमेनाइट रेखा (5782) के नीचे,
सिलिमेनाइट कायनाइट से अधिक स्थायी है, और ऐंडालुसाइट $=$ सिलिमेनाइट रेखा
(1799) के ऊपर, ऐंडालुसाइट से अधिक स्थायी है: बिंदु सिलिमेनाइट क्षेत्र में स्थित है, और वहाँ
कायनाइट $=$ ऐंडालुसाइट साम्य मितस्थायी है। (ग) \qty{4}{kbar} पर: 1799
और \qty{5782}{bar} के बीच, सिलिमेनाइट। \qty{7}{kbar} पर: \qty{5782}{bar} के ऊपर कायनाइट
सिलिमेनाइट से अधिक स्थायी है, और \qty{4885}{bar} के ऊपर वह ऐंडालुसाइट से भी: कायनाइट।
\end{solution}

\begin{solution}{exo:b2:metamorphic-facies:11}
(क) $z = P/(\rho g) = 3.75\times10^8/(2700\times9.81) = \qty{14.2}{km}$, अनिश्चितता
$0.25\times10^8/26487 = \qty{0.94}{km}$। (ख) $(504 - 10)/14.16 =
\qty{35}{\celsius.km^{-1}}$। (ग) तीनों खनिज केवल \omterm{def:b2:metamorphic-facies:invariant}{निश्चर बिंदु} पर साम्य में
सहअस्तित्व में रहते हैं: चट्टान का पथ उससे होकर जाना चाहिए, और अभिक्रियाएँ
वहीं रुकनी चाहिए। अधिकांश चट्टानें उसके पास से गुज़रती हैं और एक या दो बहुरूप रखती हैं।
\end{solution}

\begin{solution}{exo:b2:metamorphic-facies:12}
(क) X: $T = 10 + 10\times30 = \qty{310}{\celsius}$ (\qty{583.15}{K}); $P = 3000\times
9.81\times30000 = \qty{8.83e8}{Pa} = \qty{8.83}{kbar}$। Y: $T = 10 + 30\times20 =
\qty{610}{\celsius}$ (\qty{883.15}{K}); $P = 2800\times9.81\times20000 = \qty{5.49}{kbar}$।
(ख) X: कैल्साइट $=$ ऐरागोनाइट $13.31\times583.15 = \qty{7.76}{kbar} < 8.83$ पर, ऐरागोनाइट;
एल्बाइट $=$ जेडाइट + क्वार्ट्ज़ $(32.4\times583.15 - 5000)/1.698 = \qty{8.18}{kbar} <
8.83$ पर, जेडाइट + क्वार्ट्ज़; कायनाइट $=$ ऐंडालुसाइट $(8.59\times583.15 - 3900)/0.737 =
\qty{1.51}{kbar}$ पर, त्रिक बिंदु के तापमान से नीचे: कायनाइट। Y: कैल्साइट $=$
ऐरागोनाइट \qty{11.75}{kbar} $> 5.49$ पर, कैल्साइट; एल्बाइट $=$ जेडाइट + क्वार्ट्ज़
\qty{13.9}{kbar} पर, एल्बाइट; त्रिक बिंदु से ऊपर, कायनाइट $=$ सिलिमेनाइट
\qty{6.00}{kbar} $> 5.49$ पर और ऐंडालुसाइट $=$ सिलिमेनाइट \qty{1.56}{kbar} $< 5.49$ पर:
सिलिमेनाइट। (ग) X, \qty{10}{\celsius\per\km} पर, अधिक ठंडा पथ है: केवल
\qty{310}{\celsius} पर जेडाइट और ऐरागोनाइट उसे नीलशिस्ट संलक्षणी में रखते हैं। Y, तीन गुना
अधिक गर्म, \qty{610}{\celsius} और मध्यम दाब पर सिलिमेनाइट के साथ, 
एम्फ़ीबोलाइट संलक्षणी में समाप्त होता है।
\end{solution}

\begin{solution}{pb:b2:metamorphic-facies:1}
\noindent\textbf{भाग I।}

\textbf{1.} $\Delta_r H = -2589.9 + 2593.8 = \qty{3.9}{kJ}$; $\Delta_r S = 91.39 - 82.80 =
\qty{8.59}{J.K^{-1}}$; $\Delta_r V = 5.152 - 4.415 = \qty{0.737}{J.bar^{-1}}$।

\textbf{2.} $\qty{7.7}{kJ}$; $\qty{12.60}{J.K^{-1}}$; $\qty{0.571}{J.bar^{-1}}$।

\textbf{3.} $\qty{3.8}{kJ}$; $\qty{4.01}{J.K^{-1}}$; $\qty{-0.166}{J.bar^{-1}}$; हर एक
दूसरे में से पहले को घटाने पर आता है, क्योंकि ऐंडालुसाइट कट जाता है।

\textbf{4.} सिलिमेनाइट की एन्ट्रॉपी सबसे अधिक है (95.40): वह अपनी दोनों रेखाओं के उच्च-तापमान
पक्ष पर स्थित है। कायनाइट का आयतन सबसे कम है (4.415): वह अपनी दोनों रेखाओं के
उच्च-दाब पक्ष पर स्थित है। ऐंडालुसाइट, सबसे बड़ा आयतन, निम्न-दाब
पक्ष रखता है।

\textbf{5.} $8.59/0.737 = \qty{11.66}{bar/K}$; $12.60/0.571 = \qty{22.07}{bar/K}$;
$4.01/(-0.166) = \qty{-24.16}{bar/K}$।

\textbf{6.} सिलिमेनाइट में ऐंडालुसाइट से अधिक एन्ट्रॉपी पर कम आयतन है: $\Delta_r S > 0$
और $\Delta_r V < 0$। तापमान बढ़ाना सिलिमेनाइट के पक्ष में है, दाब बढ़ाना
भी, इसलिए रेखा अधिक तापमान पर कम दाब की ओर जाती है।

\noindent\textbf{भाग II।}

\textbf{7.} $3900 - 298.15\times8.59 = \qty{+1339}{J}$; $7700 - 298.15\times12.60 =
\qty{+3943}{J}$। दोनों धनात्मक: सतह पर कायनाइट स्थायी है।

\textbf{8.} $T = \Delta_r H/\Delta_r S$: $3900/8.59 = \qty{454.0}{K} = \qty{181}{\celsius}$;
$7700/12.60 = \qty{611.1}{K} = \qty{338}{\celsius}$; $3800/4.01 = \qty{947.6}{K} =
\qty{674}{\celsius}$।

\textbf{9.} \qty{1}{bar} पर गर्म करने पर कायनाइट \qty{181}{\celsius} पर ऐंडालुसाइट को
स्थान देता है; \qty{338}{\celsius} पर कायनाइट अब स्थायी नहीं है, इसलिए वहाँ कायनाइट $=$
सिलिमेनाइट मितस्थायी है; ऐंडालुसाइट \qty{674}{\celsius} पर सिलिमेनाइट को स्थान
देता है। पहला और तीसरा स्थायी साम्य हैं।

\textbf{10.} $(8.59\times673.15 - 3900)/0.737 = \qty{2554}{bar} = \qty{2.55}{kbar}$,
प्रयोग से \qty{0.01}{kbar} दूर।

\textbf{11.} $(12.60\times973.15 - 7700)/0.571 = 4562/0.571 = \qty{7989}{bar} =
\qty{7.99}{kbar}$, प्रयोग से \qty{0.13}{kbar} ऊपर।

\textbf{12.} कायनाइट $=$ ऐंडालुसाइट \qty{1}{bar} पर \qty{181}{\celsius} से उठती है, ढलान
\qty{11.7}{bar/K} के साथ; ऐंडालुसाइट $=$ सिलिमेनाइट \qty{200}{\celsius} पर लगभग \qty{11.5}{kbar} से
गिरकर \qty{674}{\celsius} पर \qty{1}{bar} तक आती है; कायनाइट $=$ सिलिमेनाइट, दोनों उठती रेखाओं में
अधिक खड़ी, दोनों को लगभग \qty{514}{\celsius} और \qty{3.9}{kbar} पर काटती है, और हर रेखा प्रतिच्छेद के केवल
एक ओर स्थायी है (अध्याय का आरेख)।

\noindent\textbf{भाग III।}

\textbf{13.} $a_1 = 8.59/0.737 = \qty{11.655}{bar/K}$, $b_1 = 3900/0.737 =
\qty{5291.7}{bar}$; $a_2 = 12.60/0.571 = \qty{22.067}{bar/K}$, $b_2 = 7700/0.571 =
\qty{13485.1}{bar}$।

\textbf{14.} $T = (b_2 - b_1)/(a_2 - a_1) = 8193.4/10.412 = \qty{786.92}{K} =
\qty{513.8}{\celsius}$।

\textbf{15.} $P = 11.655\times786.92 - 5291.7 = \qty{3880}{bar} = \qty{3.88}{kbar}$।

\textbf{16.} $(4.01\times786.92 - 3800)/(-0.166) = (3155.5 - 3800)/(-0.166) =
\qty{3882}{bar}$: आँकड़ों के पूर्णांकन की सीमा तक वही दाब, जैसा अध्याय की प्रतिज्ञप्ति माँगती है।

\textbf{17.} \qty{514}{\celsius}, \qty{504}{\celsius} से \qty{10}{\celsius} दूर है, यानी
\qty{20}{\celsius} की अनिश्चितता के भीतर; \qty{3.88}{kbar}, 3.75 से \qty{0.13}{kbar} दूर है,
यानी 0.25 के भीतर। सहमति अच्छी है।

\textbf{18.} नहीं: \qty{536}{\celsius} प्रयोग से \qty{32}{\celsius} दूर है, उसकी
अनिश्चितता के बाहर। ऊष्मा धारिताओं की उपेक्षा मुख्य त्रुटि नहीं थी: एन्थैल्पियाँ हैं।

\noindent\textbf{भाग IV।}

\textbf{19.} $\Delta_r H$, कायनाइट $\to$ ऐंडालुसाइट का, \qty{1000}{J} बढ़ता है; स्थिर
$T$ पर, $P = (T\Delta_r S - \Delta_r H)/\Delta_r V$ में $1000/0.737 = \qty{1357}{bar}$ की कमी आती है।

\textbf{20.} $\Delta T = -1357/(a_2 - a_1) = -1357/10.41 = \qty{-130}{K}$: त्रिक बिंदु
लगभग \qty{384}{\celsius} पर खिसक जाता है।

\textbf{21.} कायनाइट $=$ सिलिमेनाइट रेखा $1000/0.571 = \qty{1751}{bar}$ नीचे आती है, और
त्रिक बिंदु $+1751/10.41 = \qty{+168}{K}$ खिसकता है।

\textbf{22.} हर एन्थैल्पी अकेले, $\pm\qty{2.0}{kJ}$ बदलने पर, त्रिक बिंदु को
$\pm\qty{260}{K}$ (ऐंडालुसाइट) या $\pm\qty{336}{K}$ (सिलिमेनाइट) खिसकाती है: अकेली कैलोरीमिति
कई सौ अंशों में फैले त्रिक बिंदुओं की अनुमति देती है।

\textbf{23.} सारणीबद्ध एन्थैल्पियाँ पूरी तरह कैलोरीमितीय नहीं हैं: अपनी
अनिश्चितताओं के भीतर उन्हें इस प्रकार समायोजित किया जाता है कि परिकलित रेखाएँ उत्क्रमित प्रयोगों से मेल खाएँ।
प्रयोग आरेख को स्थिर करते हैं; तालिकाएँ उसे, संगत रूप से, अन्य अभिक्रियाओं तक ले जाती हैं।

\textbf{24.} स्थिर गुणों के साथ तालिकाओं से परिकलित त्रिक बिंदु
\textbf{\qty[text-series-to-math=true]{514}{\celsius}} पर स्थित है।
\end{solution}
